
\section{固有态射}

拓扑空间之间的映射（也就是\trContinuousMap！）称为\textbf{固有的}（proper），若任何\trCompact{}集合的原像都是\trCompact{}的。\trMorphism{}的\trProperness{}是一条与之类似的性质。例如，$\mathbb{C}$ 上的一个\trVariety{}若在\trClassicalTopology{}下\trCompact{}，就会是固有的（参看 [Se3, \S7]）。另一方面，我们将会看到：射影 $A$-\trScheme{}都在 $A$ 上固有——所以这是一条射影概形都满足的良好性质，而且用起来也很顺手。

回想（\S8.3.5）：拓扑空间之间的（\trContinuous）映射 $\pi\colon X\to Y$ 称为\emph{闭}的，若对每个闭子集 $S\subset X$，$\pi(S)$ 也是闭的。概形之间的\trMorphism{}称为闭的，若其底下的\trContinuousMap{}是闭的。我们称概形间的\trMorphism{} $\pi\colon X\to Y$ 是\trUniversallyClosed{}的，若对每个\trMorphism{} $Z\to Y$，诱导的\trMorphism{} $Z\times_Y X\to Z$ 都是闭的。换言之，一个\trMorphism{}是泛闭的，若它在任何\trBaseChange{}之下都保持为闭的。事到如今，你大概会自动期待下面这道练习了。

\textbf{11.4.A.（容易的练习）}\quad 证明：泛闭\trMorphism{}构成的类是 \S8.1 意义下的「合理」类。

为给概形的\trProperness{}的定义提供动机，我们指出：设 $\pi\colon X\to Y$ 是局部紧\trHausdorff{}空间之间的一个\trContinuousMap，且二者的拓扑都具有可数\trBase，那么 $\pi$ 是泛闭的，当且仅当它是固有的（即\trCompact{}子集的原像是\trCompact{}的）。欢迎你把这件事当作一道练习证一证。

\textbf{11.4.1.}\ \emph{\trhDefinition{}。}\quad \trMorphism{} $\pi\colon X\to Y$ 称为\textbf{\trProper}的，若它既是\trSeparated{}的、\trFiniteType{}的，又是\trUniversallyClosed{}的。一个\trScheme{} $X$ 常被说成是固有的，若它某个隐含的\trStructureMorphism{}是固有的。例如，一个 $k$-\trScheme{} $X$ 常被描述为固有的，若 $X\to\Spec k$ 是固有的。（固有的 $k$-\trScheme{}也常说成是\textbf{\trComplete}的（complete）。我们不会使用这一术语。）若 $A$ 是一个\trRing，人们常说一个 $A$-\trScheme{}是\textbf{在 $A$ 上固有}的，若它在 $\Spec A$ 上固有。

让我们在实践中把这条观念试用一下。我们预期 $\mathbb{A}^1_{\mathbb{C}}\to\Spec\mathbb{C}$ 不是固有的，因为与 $\mathbb{A}^1_{\mathbb{C}}$ 相应的\trComplex\trManifold{}不是\trCompact{}的。然而请注意：这个映射是\trSeparated{}的（它是仿射概形之间的映射）、\trFiniteType{}的，而且是（平凡地）闭的。所以，关键全在「泛」字上。

\textbf{11.4.B.（容易的练习）}\quad 证明：$\pi\colon\mathbb{A}^1_{\mathbb{C}}\to\Spec\mathbb{C}$ 是一个\trClosed{}映射，但不是泛闭的。从而 $\pi$ 不是固有的。可能的提示：考虑一个精心挑选的映射，比如 $\mathbb{A}^1_{\mathbb{C}}\times_{\mathbb{C}}\mathbb{P}^1_{\mathbb{C}}\to\mathbb{P}^1_{\mathbb{C}}$。（\mbox{图~19.1}给出了另一个有限型、分离、闭、却不固有的\trMorphism{}的例子。而要证明\emph{那个}\trMorphism{}不固有，则需要更多的创造力。）

\textbf{11.4.2.}\ \emph{\trhExample{}。}\quad 作为第一个例子：闭嵌入都是固有的。它们显然\trSeparated，因为仿射\trMorphism{}是\trSeparated{}的（命题~11.2.4）。它们是\trFiniteType{}的。经过\trBaseChange{}之后，它们仍是闭嵌入（\S10.2.2），而闭嵌入总是闭的。这条可以轻易地进一步推广，如下。

\textbf{11.4.3.\ \trhProposition{}。}\ ——\ \emph{有限\trMorphism{}都是固有的。}

\emph{证明。}\ 有限\trMorphism{}是\trSeparated{}的（它们按定义是仿射的，而仿射\trMorphism{}是\trSeparated{}的，命题~11.2.4），并且是\trFiniteType{}的（基本原因：环上的有限\trModule{}自动有限生成）。为证有限\trMorphism{}在任意\trBaseChange{}之后都是闭的，我们注意：它们在任意\trBaseChange{}之后仍是有限的（有限性对\trBaseChange{}保持，\mbox{练习~10.4.D(d)}），而有限\trMorphism{}是闭的（\mbox{练习~8.3.L}）。\hfill$\square$

\textbf{11.4.4.\ \trhProposition{}。}\ ——\ \emph{(a) 「固有\trMorphism」这一观念对\trBaseChange{}稳定。}\\
\emph{(b) 「固有\trMorphism」这一观念是对目标局部的（也就是说，$\pi\colon X\to Y$ 是固有的，当且仅当对任意仿射开覆盖 $\mathsf{U}_i\to Y$，$\pi^{-1}(\mathsf{U}_i)\to\mathsf{U}_i$ 都是固有的）。注意「仅当」方向由 (a) 推出——考虑经 $\mathsf{U}_i\to Y$ 的\trBaseChange{}即可。}\\
\emph{(c) 「固有\trMorphism」这一观念对\trComposition{}保持。}\\
\emph{(d) 两个固有\trMorphism{}的积是固有的：若 $\pi\colon X\to Y$ 与 $\pi'\colon X'\to Y'$ 都是固有的（其中诸\trMorphism{}均为某 $Z$-\trScheme{}之间的\trMorphism），则 $\pi\times\pi'\colon X\times_Z X'\to Y\times_Z Y'$ 是固有的。}\\
\emph{(e) 设}
\begin{equation*}\tag{11.4.4.1}\label{eq:properness-cancellation-triangle}
\begin{tikzcd}[row sep=2.4em, column sep=3.2em]
X \ar[rr, "\pi"] \ar[dr, "\tau"'] & & Y \ar[dl, "\rho"] \\
& Z &
\end{tikzcd}
\end{equation*}
\emph{是一个\trCommutativeDiagram，其中 $\tau$ 是固有的，而 $\rho$ 是\trSeparated{}的。那么 $\pi$ 是固有的。}

(e) 的一个示范应用：从固有的 $k$-\trScheme{}出发、（在 $\Spec k$ 上）到\trSeparated{}的 $k$-\trScheme{}的\trMorphism{}总是固有的。

\emph{证明。}\ (a) 至 (c) 说的是：固有\trMorphism{}构成 \S8.1 意义下的一个「合理」的\trMorphism{}类。于是它们得以成立，因为\trSeparated{}态射、有限型态射与泛闭态射都是「合理」的类。然后，(d) 由 (a) 与 (c) 借\cref{ex:property-product-preserved}推出。闭嵌入是固有的（例~11.4.2），故 (e) 由（对固有\trMorphism{}的）\trCancellationTheorem~11.1.1 推出。\hfill$\square$

\textbf{11.4.C.（练习：「固有概形的像是固有的」）}\quad 设在图
\[
\begin{tikzcd}[row sep=2.4em, column sep=3.2em]
X \ar[rr, "\pi"] \ar[dr, "\tau"'] & & Y \ar[dl, "\rho"] \\
& Z &
\end{tikzcd}
\]
中 $\tau$ 是固有的，而 $\rho$ 是\trSeparated{}且有限型的。证明：$X$ 在 $\pi$ 之下的\trSchemeTheoreticImage{}是一个固有的 $Z$-\trScheme。（由此可见，代数几何中的固有性表现得恰如拓扑中的固有性一样。）

现在我们来看固有态射最重要的例子。

\textbf{11.4.5.\ \trhTheorem{}。}\ ——\ \emph{射影 $A$-\trScheme{}都在 $A$ 上固有。}

（由于到 $\Spec A$ 的有限\trMorphism{}都是射影 $A$-\trScheme{}（\mbox{练习~8.3.I}），定理~11.4.5 可以用来给出「有限\trMorphism{}都是固有的」（命题~11.4.3）的第二个证明。）

\emph{证明。}\ 射影 $A$-\trScheme{}的\trStructureMorphism{} $X\to\Spec A$ 可以分解为一个\trClosedEmbedding{}复合上 $\mathbb{P}^n_A\to\Spec A$。闭嵌入是固有的（例~11.4.2），而固有\trMorphism{}的\trComposition{}是固有的（命题~11.4.4），故只须证明 $\mathbb{P}^n_A\to\Spec A$ 是固有的。我们已经见过这个态射是\trFiniteType{}的（容易练习~5.3.D）且是\trSeparated{}的（命题~11.2.8），故只须证明 $\mathbb{P}^n_A\to\Spec A$ 是泛闭的。由于 $\mathbb{P}^n_A=\mathbb{P}^n_{\mathbb{Z}}\times_{\mathbb{Z}}\Spec A$，只须证明：对任何概形 $X$，$\mathbb{P}^n_X:=\mathbb{P}^n_{\mathbb{Z}}\times_{\mathbb{Z}}X\to X$ 都是闭的。但「闭」这条性质是对目标（此处即 $X$）局部，所以用仿射开子集盖住 $X$ 之后，便只须证明：对一切的环 $B$，$\mathbb{P}^n_B\to\Spec B$ 都是闭的。这正是\trFundamentalTheoremOfEliminationTheory（定理~8.4.10）。\hfill$\square$

\textbf{11.4.6.}\ \emph{\trhRemark{}：}\textbf{「合理」的固有概形都是射影的。}\quad 要举出一个固有却不射影的 $A$-\trScheme{}的例子，可不是件容易事！在域上，所有固有\trCurve{}都是射影的（见注记~19.2.I），所有光滑\trSurface{}都是射影的（\trZariski{}这条定理的证明见 [Ba, Thm.~1.28]；光滑性当然还没定义）。我们将在 \S15.2.11 遇到一个固有却不射影的\trVariety{}（一个奇异三维簇）的第一例。我们日后将在 \S19.11.11 见到一个固有非射影\trSurface{}的例子，并在\mbox{练习~20.2.G}见到一个更简单的。等我们了解了平坦性之后，还将看到\trHironaka{}那个 $\mathbb{C}$ 上固有、非射影、不可约且光滑的三维簇的例子（\S24.7.6）。

\textbf{11.4.7. 连通既约的固有 $\overline{k}$-\trScheme{}上的函数必为常数。}\quad
作为这些观念的一个颇具启发性的应用，我们来证明：若 $X$ 是一个连通、既约、固有的 $k$-\trScheme{}，其中 $k=\overline{k}$，则 $\Gamma(X,\OO_X)=k$。连通\trCompact{}\trComplex\trManifold{}的类似事实用到的是最大值原理（maximum principle）。我们在 $X=\mathbb{P}^n$ 的特殊情形见过它（\mbox{练习~4.4.E}）；它日后将被\trGrothendieck{}凝聚定理~18.9.1（Coherence Theorem）大大地推广。

设 $f\in\Gamma(X,\OO_X)$（$f$ 是 $X$ 上的一个函数）。这与一个映射 $\pi\colon X\to\mathbb{A}^1_k$ 是同一回事（\mbox{练习~7.3.G}，进一步的讨论见 \S7.6.1）。令 $\pi'$ 为 $\pi$ 复合上开嵌入 $\mathbb{A}^1\hookrightarrow\mathbb{P}^1$ 的结果。由命题~11.4.4(e)，$\pi'$ 是固有的，从而特别地是闭的。由于 $X$ 连通，$\pi'$ 的像也连通。于是 $\pi'$ 的集合论像要么是一个闭点，要么是整个 $\mathbb{P}^1$。但 $\pi'$ 的集合论像落在 $\mathbb{A}^1$ 中，故它必是一个闭点 $p$（我们把它与 $k$ 的一个元素相认同）。

由推论~9.4.5，$\pi$ 的\trSchemeTheoreticImage{}的\trSupport{}是这个闭点 $p$。由\mbox{练习~9.4.A}，\trSchemeTheoreticImage{}恰好就是 $p$（带既约结构）。于是 $\pi$ 可以解读为：先取到 $\Spec k$ 的\trStructureMorphism，再复合一个把 $\Spec k$ 认同于 $p$ 的、到 $\mathbb{A}^1$ 的闭嵌入。你应当能够验证：这正是对应于常值函数 $f=p$ 的那个到 $\mathbb{A}^1$ 的映射。

（若放松不同的假设，反例为何？条件 $k=\overline{k}$ 是不可忽视的——若 $k=\mathbb{R}$，则 $X=\Spec\mathbb{C}$ 是一个连通既约的固有射影 $k$-\trScheme，而 $\Gamma(X,\OO_X)=k^2$。一个有趣的事实：$\mathbb{P}^2_k$ 中 $x^2=0$ 表明，即便是非既约的固有 $k$-\trScheme，其上的函数也可以只有常数；见 \S5.2.3 与\mbox{练习~18.6.U}。）

\textbf{11.4.8. 或许能帮助你正确思考有限性的若干事实（尚未证明）。}\quad
下面的几条事实或许能照亮「有限性」这个观念。我们日后再证明它们。

\trLocally\trNoetherian{}概形之间的\trMorphism{}是有限的，当且仅当它既固有又仿射，当且仅当它既固有又拟有限（定理~28.5.2）。这个陈述的若干部分我们已经证过，但只有等到我们知道\trZariskiMainTheorem{}之后，才能完成整个证明；对比 \S8.3.12。

作为一个应用：从固有概形到分离概形的拟有限\trMorphism{}是有限的。道理如下：设 $\pi\colon X\to Y$ 是 $Z$ 上的一个拟有限\trMorphism，其中 $X$ 在 $Z$ 上固有。那么由（对固有\trMorphism{}的）\trCancellationTheorem~11.1.1，$X\to Y$ 是固有的。于是既然 $\pi$ 既拟有限又固有，$\pi$ 就是有限的。

作为一个显式的例子，考虑由 $[x,y]\mapsto[f(x,y),\,g(x,y)]$ 给出的映射 $\pi\colon\mathbb{P}^1_k\to\mathbb{P}^1_k$，其中 $f$ 与 $g$ 是次数相同、且在 $\mathbb{P}^1$ 中没有公共根的齐次多项式。诸\trFiber{}有限，而 $\pi$ 是固有的（由对固有\trMorphism{}的\trCancellationTheorem~11.1.1，如命题~11.4.4\footnote{原书此处作（Theorem 11.4.4），而 11.4.4 实为命题，当系「命题」之误，径改。——译者注} 之后所讨论的那样），故 $\pi$ 是有限的。这件事固然也可以直接验证，但现在我们可以省去这份烦冗了。

\textbf{11.4.9.}\ \extrahard\ \textbf{\trGroupVariety{}。}\quad \S11.4 的其余部分都是双星的，而且我们也将通盘在域 $k$ 上的\emph{\trVariety}这一\trCategory{}中工作。我们简短地讨论\trGroupVariety{}，主要是因为力所能及，而不是因为我们有什么深刻的见地要说。之所以此刻就把它们拿出来讨论，只是因为固有性给出一条出人意料的推论：固有的\trGroupVariety{}是\trAbelian{}的（多亏出奇简单的刚性引理~11.4.12）。

如 \S11.2.9 所讨论的，既然我们手里已有域 $k$ 上的簇这一\trCategory{}，便立即有了\textbf{\trGroupVariety{}}（group variety）的观念（多谢 \S7.6.4 讲的任意\trCategory{}中的\trGroupObject）。一个（$k$ 上的）\textbf{\trAlgebraic\trGroup}（algebraic group）是一个光滑的\trGroupVariety{}。例子包括 $\mathrm{GL}_n$（\mbox{练习~7.6.N(b)}）——$\mathbb{G}_m$ 是其特例——以及 $\mathrm{SL}_n$（\mbox{练习~7.6.Q}）。

\textbf{11.4.10.}\ \emph{旁注}（我们既不证明也不使用）。\quad \trGroupVariety{}在特征 $0$ 时自动光滑（从而是\trAlgebraicGroup{}）。\trGroupVariety{}在正特征下未必光滑。（设 $k$ 是特征 $p$ 的不\trPerfectField，取 $t\in k$ 使其在 $k$ 中没有 $p$ 次根。考虑「加法群」（additive group）$\mathbb{G}_a^2=\Spec k[x,y]$ 中由 $x^p=ty^p$ 给出的「闭子群概形」。这个群概形既约，却不几何既约；而且不光滑，甚至连正则都不是。）\trAlgebraicGroup{}自动拟射影。\trAlgebraicGroup{} $G$ 是仿射的（作为概形），当且仅当它能容纳一个到某 $\mathrm{GL}_n$（$n\ge 0$）的\trClosedImmersion{}，使得 $G\to\mathrm{GL}_n$ 是群概形之间的同态（\mbox{练习~7.6.N(a)}）。关于这些问题与其他议题的进一步讨论，见 [P2, \S5]。

\textbf{11.4.11.}\ \emph{\trhDefinition{}：}\textbf{\trAbelianVariety{}}。\quad 我们现在可以定义\trVariety{}最重要的类之一了：\trAbelianVariety{}（abelian variety）。坏消息是，我们还没有一个真正的例子。我们日后将见到椭圆曲线作为一例（\S19.10），但比它更奇特的例子我们就无缘得见了。尽管如此，能这么早就作出这个定义，还是令人愉快的。

域 $k$ 上的一个\textbf{\trAbelian\trVariety}（abelian variety）是一个几何整且（在 $k$ 上）射影的\trAlgebraicGroup{}。事实证明，$\mathbb{C}$ 上任何\trAbelianVariety{}的\trAnalytification{}都是一个\trComplexTorus{}（complex torus），也就是说形如 $\mathbb{C}^n/\Lambda$，其中 $\Lambda$ 是一个\trLattice{}（lattice）（秩为 $2n$）。（关键想法：连通\trCompact{}\trComplex{}\trLie{}群是交换的；见注记~11.4.14；其\trUniversalCover{}（universal cover）是一个\trSimplyConnected{}（simply connected）的交换复\trLie{}群，从而是 $\mathbb{C}^n$。参看 [BL, Lem.~1.1.1]。）

我们现在来证明：\trAbelianVariety{}是\trAbelian{}的，也就是说，\trAbelianVariety{}是交换的代数\trGroupVariety{}。（这话听起来像同义反复，其实不然！形容词「阿贝尔」用于两种不同的意义；二者的共同点是：它们都源自 Niels Henrik Abel（尼尔斯·亨利克·阿贝尔）之名。）由于\trAlgebraicGroup{}未必交换（$\mathrm{GL}_n$ 便是见证），交换性竟能被「紧致性」逼出来，这就多少令人称奇了。

关键的结果是下面这个事实，它本身就美得很。

\textbf{11.4.12.\ 刚性引理}（Rigidity Lemma）\textbf{。}\ ——\ \emph{设 $X$、$Y$、$Z$ 都是\trVariety，其中 $X$ 固有且\trGeometricallyIntegral，并带有一个 $k$-\trValuedPoint{} $p$，而 $Y$ 不可约。设 $\mathrm{pr}\colon X\times Y\to Y$ 是\trProjection，而 $\alpha\colon X\times Y\to Z$ 是任一\trMorphism。设 $q\in Y$ 使得 $\alpha$ 在 $X\times\{q\}$ 上为常值；比如说 $\alpha(X\times\{q\})=r\in Z$。}\\
\emph{(a) 那么存在\trMorphism{} $\psi\colon Y\to Z$ 使 $\alpha=\psi\circ\mathrm{pr}$。（特别地，对每个 $q'\in Y$，$\alpha$ 在 $X\times\{q'\}$ 上为常值。）}\\
\emph{(b) 若 $\alpha$ 在 $\{p\}\times Y$ 上也为常值，则 $\alpha$ 在整个 $X\times Y$ 上都是常值的。}

\[
\begin{tikzcd}[row sep=3em, column sep=3.6em]
& & Z \\
X\times\{q\} \ar[r] \ar[d] & X\times Y \ar[ur, "\alpha"] \ar[d, "\mathrm{pr}"', shift left=0.55em] \\
q \ar[r] & Y \ar[u, "\sigma_p"', shift right=0.55em] \ar[uur, bend right=24, "\psi"]
\end{tikzcd}
\]

这些假设可以在许多方向上放松，但这个版本对我们的目的已经够用。如果你没有读过讨论几何整性的 \S10.5.1，也不必过于担心；当 $k=\overline{k}$ 且 $X$ 是整的时候，它是自动成立的（「更难事实」~10.5.3）。即便你读过 \S10.5.1，先在 $k=\overline{k}$ 的假定下读这个证明也会很有帮助，以免被几何点岔开对主要想法的注意力。但话又说回来，\S11.4.9 的全部内容都是双星的，所以你本来就不该读到这里的。

\textbf{11.4.D.（练习）}\quad 证明：$X$ 的固有性假设是必要的。（你能让这与你的直觉相安无事吗？）

\emph{刚性引理~11.4.12 的证明。}\ 定义 $\beta\colon X\times Y\to Z$ 为 $\beta(x,y)=\alpha(p,y)$。（更精确地说，若 $\sigma_p$ 是从 $\Spec k\hookrightarrow p$ 拉回的 $X\times Y\to Y$ 的\trSection，则 $\beta=\alpha\circ\sigma_p\circ\mathrm{pr}$。）我们将证明 $\beta(x,y)=\alpha(x,y)$。这将蕴涵 (a)（取 $\psi=\alpha\circ\sigma_p$），而 (a) 又立即蕴涵 (b)。

命题~10.5.19（几何既约乘既约是既约）蕴涵 $X\times Y$ 是既约的。\mbox{练习~10.5.O}（几何不可约乘不可约是不可约）蕴涵 $X\times Y$ 是不可约的。由于 $Z$ 是分离的，由\mbox{练习~11.3.A}，只须证明 $\alpha=\beta$ 在 $X\times Y$ 的某个非空（从而稠密）开子集上成立。

设 $\mathsf{U}\subset Z$ 是 $r$ 的一个仿射开邻域。那么 $\alpha^{-1}(Z\setminus\mathsf{U})$ 是 $X\times Y$ 的闭子集。由于\trProjection $\mathrm{pr}$ 是固有的（利用 $X$ 的固有性，以及固有性在\trBaseChange{}下的保持，命题~11.4.4(a)），$\mathrm{pr}(\alpha^{-1}(Z\setminus\mathsf{U}))$ 是闭的。它的补集是开的，并且含有 $q$；令 $\mathsf{V}\subset Y$ 是 $q$ 的一个开邻域，与 $\mathrm{pr}(\alpha^{-1}(Z\setminus\mathsf{U}))$ 不相交。

\textbf{11.4.E.（练习）}\quad 设 $\gamma\colon\Spec\overline{K}\to\mathsf{V}$ 是 $\mathsf{V}$ 的一个\trGeometricPoint。证明：$\Spec\overline{K}$ 之上的纤维 $X_{\overline{K}}$（或者更精确地说，$\gamma$ 之上的纤维）被 $\alpha$ 限制在 $X_{\overline{K}}$ 上的 $\alpha_{\overline{K}}$ 送到 $Z$ 中的一个点。提示：重温 \S11.4.7。
\[
\begin{tikzcd}[row sep=3em, column sep=3.8em]
& & Z \\
X_{\overline{K}} \ar[rr] \ar[d] \ar[urr, "\alpha_{\overline{K}}", shift left=0.5em] & & X\times\mathsf{V} \ar[u, "\beta_{\overline{K}}", shift left=0.5em] \ar[u, bend right=30, "\alpha|_{\mathsf{V}}"'] \ar[d, "\mathrm{pr}"'] \\
\Spec\overline{K} \ar[rr, "\gamma"'] & & \mathsf{V}
\end{tikzcd}
\]
由\mbox{练习~11.3.B}（代数闭域上的簇由闭点上的映射决定），$\alpha_{\overline{K}}=\beta_{\overline{K}}$。

\textbf{11.4.F.（练习）}\quad 用\mbox{练习~11.3.A}（$\alpha=\beta$ 之处是一个\trClosedSubscheme{}这一事实——而且它表示某个函子）证明：在 $\mathrm{pr}^{-1}(\mathsf{V})$ 上有 $\alpha=\beta$。\hfill$\square$

\textbf{11.4.13.\ \trhCorollary{}。}\ ——\ \emph{设 $A$ 是一个\trAbelian\trVariety，$G$ 是一个\trGroup\trVariety，而 $\phi\colon A\to G$ 是\trVariety{}之间的任一\trMorphism。那么 $\phi$ 是一个\trTranslation{}（translation）与一个同态的复合。}

\emph{证明。}\ 设 $\mathrm{e}_A$ 与 $\mathrm{e}_G$ 分别是 $A$ 与 $G$ 的\trIdentityPoint{}（identity point）。把 $\phi$ 复合上一个\trTranslation{}，便可设 $\phi(\mathrm{e}_A)=\mathrm{e}_G$。考虑由 $\alpha(a_1,a_2)=\phi(a_1a_2)\phi(a_2)^{-1}\phi(a_1)^{-1}$ 给出的\trMorphism{} $\alpha\colon A\times A\to G$。那么 $\alpha(\{\mathrm{e}_A\}\times A)=\alpha(A\times\{\mathrm{e}_A\})=\{\mathrm{e}_G\}$，故由刚性引理~11.4.12，$\alpha$ 是常值的，并把 $A\times A$ 送到 $\mathrm{e}_G$。于是 $\phi(a_1a_2)=\phi(a_1)\phi(a_2)$，故 $\phi$ 是一个同态。\hfill$\square$

\textbf{11.4.G.（练习：\trAbelianVariety{}是\trAbelian{}的）}\quad 证明：\trAbelianVariety{}是一个\trAbelian{}的\trGroupVariety{}。提示：把推论~11.4.13 应用到求逆\trMorphism{} $i\colon A\to A$ 上。

\textbf{11.4.14.}\ \emph{\trhRemark{}。}\quad 类似的想法还被用来证明：连通\trCompact{}\trComplex{}\trLie{}群是\trAbelian{}的；例如参看 [BL, Lem.~1.1.1]。
